\documentclass[12pt]{article}
\usepackage[utf8]{inputenc}
\usepackage[T1]{fontenc}
\usepackage{hyperref}
\title{Maiani Language Resources for K-12 Curriculum Development}
\author{Compiled by Claude Sonnet 5 (claude-sonnet-5) for Andr\'as Kornai}
\date{}
\begin{document}
\maketitle

\section{Introduction}
Maiani (ISO 639-3 tnh; Glottocode maia1253) is a Papuan language of Madang Province, Papua New Guinea. Glottolog places it in the Kaukombaran branch of the North Adelbert Range group (Nuclear Trans New Guinea). It is the southern member of a pair with Miani (code pla). Sources disagree on whether the two are separate languages or dialects. Speaker estimates range from about 2,500 to about 3,000. Searches found very little material written for schools. This report lists only what the searches returned.

\section{Primary Educational Resources}
No primers, textbooks or school programmes for Maiani turned up in the searches. If any exist, they were not found.

\section{Linguistic and Reference Materials}
\begin{itemize}
\item Wikipedia, Miani language (covers the Miani and Maiani varieties): \url{https://en.wikipedia.org/wiki/Miani_language}
\item Grambank entry for Maiani: \url{https://grambank.clld.org/languages/maia1253}
\item CrossGram entry for Maiani: \url{https://crossgram.clld.org/languages/maia1253}
\item ASJP word-list entry: \url{https://asjp.clld.org/languages/MAIANA}
\item Max Planck EVA, numerals of Maiani: \url{https://lingweb.eva.mpg.de/channumerals/Maiani.htm}
\item TransNewGuinea.org database (word pages): \url{https://transnewguinea.org/word/my?page=5}
\end{itemize}
Secondary sources mention a 1982 grammar by Loeweke and May covering Maiani, Miani and Mala. I did not find a link to it, so it is not listed as a resource.

\section{Religious and Community Resources}
No scripture translation or community organization was found in the searches.

\section{Digital and Audio Resources}
\begin{itemize}
\item Global Recordings Network, Maiani (tnh): \url{https://globalrecordings.net/en/language/tnh}
\item Global Recordings Network, Miani (pla), the related northern variety: \url{https://globalrecordings.net/en/language/pla}
\end{itemize}

\section{Curriculum Development Context}
Maiani is used by all adults but not all young people, and it is not known to be taught in schools. A curriculum would need new materials built from scratch. The sources above could serve as a starting point: the word lists, the numerals, the audio recordings and the 1982 grammar if it can be located. Community consultation is essential. No named local contacts were found.

\end{document}